\chapter{Maximum Principles}

In this chapter, we discuss maximum principles and their applications. Two
classes of maximum principles will be discussed, one due to Hopf and the other
to Alexandroff. The former gives estimates of solutions in terms of the
$L^\infty$-norm of the nonhomogenous terms while the latter gives the estimates in terms of the
$L^n$-norm. Applications include various a priori estimates and the moving plane
method.

\section*{2.1. Strong Maximum Principle}

Suppose $\Omega$ is a bounded and connected domain in $\mathbb{R}^n$. Consider the operator
$L$ in $\Omega$
\[
Lu \equiv a_{ij}(x)D_{ij}u + b_i(x)D_iu + c(x)u,
\]
for $u \in C^2(\Omega) \cap C(\overline{\Omega})$. We always assume that $a_{ij}$, $b_i$ and $c$ are continuous and
hence bounded in $\overline{\Omega}$ and that $L$ is uniformly elliptic in $\Omega$ in the following sense
\[
a_{ij}(x)\xi_i\xi_j \ge \lambda |\xi|^2 \qquad \text{for any } x \in \Omega \text{ and any } \xi \in \mathbb{R}^n,
\]
for some positive constant $\lambda$.

\begin{lemma}\label{hanlin:maxprin-lemma-2-1}\dcref{Lemma 2.1}\group{ch02-maximum-principles}\source{han-lin-lecture-notes:page-0029}
Suppose $u \in C^2(\Omega)\cap C(\overline{\Omega})$ satisfies $Lu > 0$ in $\Omega$ with $c(x) \le 0$ in
$\Omega$. If $u$ has a nonnegative maximum in $\overline{\Omega}$, then $u$ cannot attain this maximum in
$\Omega$.
\end{lemma}
\begin{proof}
Suppose $u$ attains its nonnegative maximum of $\overline{\Omega}$ at $x_0 \in \Omega$. Then
$D_iu(x_0)=0$ and the matrix $(D_{ij}(x_0))$ is semi-negative definite. By the ellip-
ticity condition, the matrix $(a_{ij}(x_0))$ is positive definite. This implies $Lu(x_0) =
a_{ij}(x_0)D_{ij}u(x_0) \le 0$, which is a contradiction. $\square$
\end{proof}

\begin{remark}\label{hanlin:maxprin-remark-2-2}\dcref{Remark 2.2}\group{ch02-maximum-principles}\source{han-lin-lecture-notes:page-0029}
If $c(x) \equiv 0$, then the requirement for the nonnegativeness on $u$
can be removed. This remark holds for many results in the rest of this section.
\end{remark}

The next result is referred to as the weak maximum principle.

\begin{lemma}\label{hanlin:maxprin-lemma-2-3}\dcref{Lemma 2.3}\group{ch02-maximum-principles}\source{han-lin-lecture-notes:page-0029}
Suppose u $\in C^2(\Omega)\cap C(\overline{\Omega})$ satisfies $Lu \ge 0$ in $\Omega$ with $c(x) \le 0$ in
$\Omega$. Then $u$ attains on $\partial \Omega$ its nonnegative maximum in $\overline{\Omega}$.
\end{lemma}
\begin{proof}
For any $\epsilon > 0$, consider $w(x) = u(x) + \epsilon e^{\alpha x_1}$ with $\alpha$ to be determined.
Then we have
\[
Lw = Lu + \epsilon e^{\alpha x_1}(a_{11}\alpha^2 + b_1\alpha + c).
\]
Since $b_1$ and $c$ are bounded and $a_{11}(x) \ge \lambda > 0$ for any $x \in \Omega$, by choosing $\alpha > 0$
large enough we get
\[
a_{11}(x)\alpha^2 + b_1(x)\alpha + c(x) > 0 \qquad \text{for any } x \in \Omega.
\]
This implies $Lw > 0$ in $\Omega$. By Lemma 2.1, $w$ attains its nonnegative maximum
only on $\partial \Omega$, i.e.,
\[
\sup_{\Omega} w \le \sup_{\partial \Omega} w^+.
\]
Then we obtain
\[
\sup_{\Omega} u \le \sup_{\Omega} w \le \sup_{\partial \Omega} w^+
\le \sup_{\partial \Omega} u^+ + \epsilon \sup_{x \in \partial \Omega} e^{\alpha x_1}.
\]
We finish the proof by letting $\epsilon \to 0$.
\end{proof}

As an application, we have the uniqueness of the solution $u \in C^2(\Omega)\cap C(\overline{\Omega})$ of
the following Dirichlet boundary value problem for $f \in C(\Omega)$ and $\varphi \in C(\partial \Omega)$
\[
Lu = f \qquad \text{in } \Omega
\]
\[
u = \varphi \qquad \text{on } \partial \Omega.
\]
if $c(x) \le 0$ in $\Omega$.

\begin{remark}\label{hanlin:maxprin-remark-2-4}\dcref{Remark 2.4}\group{ch02-maximum-principles}\source{han-lin-lecture-notes:page-0030}
The boundedness of the domain $\Omega$ is essential, since it guarantees
the existence of maximum and minimum of $u$ in $\overline{\Omega}$. The uniqueness does not hold
if the domain is unbounded. Some examples are given in Section 1.1. Equally
important is the nonpositiveness of the coefficient $c$.
\end{remark}

\begin{example}\label{hanlin:maxprin-example-2-5}\dcref{Example 2.5}\group{ch02-maximum-principles}\source{han-lin-lecture-notes:page-0030}
Set $\Omega = \{(x,y) \in \mathbb{R}^2; 0 < x < \pi, 0 < y < \pi\}$. Then $u =
\sin x \sin y$ is a nontrivial solution of the problem
\[
\Delta u + 2u = 0 \qquad \text{in } \Omega
\]
\[
u = 0 \qquad \text{on } \partial \Omega.
\]
\end{example}

The next result is called the Hopf lemma.

\begin{theorem}\label{hanlin:maxprin-theorem-2-6}\dcref{Theorem 2.6}\group{ch02-maximum-principles}\source{han-lin-lecture-notes:page-0030,han-lin-lecture-notes:page-0031}
Let $B$ be an open ball in $\mathbb{R}^n$ with $x_0 \in \partial B$. Suppose $u \in$
$C^2(B)\cap C(B\cup\{x_0\})$ satisfies $Lu \ge 0$ in $B$ with $c(x) \le 0$ in $B$. Assume in addition
that
\[
u(x) < u(x_0) \qquad \text{for any } x \in B \text{ and } u(x_0) \ge 0.
\]
Then for each outward direction $v$ at $x_0$ with $v \cdot n(x_0) > 0$ there holds
\[
\liminf_{t \to 0^+} \frac{1}{t}[u(x_0) - u(x_0 - tv)] > 0.
\]
\end{theorem}
\begin{remark}\label{hanlin:maxprin-remark-2-7}\dcref{Remark 2.7}\group{ch02-maximum-principles}\source{han-lin-lecture-notes:page-0030,han-lin-lecture-notes:page-0031}
If, in addition, $u \in C^1(B\cup\{x_0\})$, we have
\[
\frac{\partial u}{\partial \nu}(x_0) > 0.
\]
\end{remark}
\begin{proof}
We assume that $B$ is centered at the origin with radius $r$. We assume
further that $u \in C(\overline{B})$ and $u(x) < u(x_0)$ for any $x \in \overline{B}\setminus\{x_0\}$, since we can construct
a ball $B_* \subset B$ tangent to $B$ at $x_0$.

Consider $v(x) = u(x) + \epsilon h(x)$ for some nonnegative function $h$. We will choose
$\epsilon > 0$ appropriately such that $v$ attains its nonnegative maximum only at $x_0$. Set
$\Sigma = B \cap B_{r/2}(x_0)$ and $h(x) = e^{-\alpha |x|^2} - e^{-\alpha r^2}$ with $\alpha$ to be determined. We check
in the following that
\[
Lh > 0 \text{ in } \Sigma.
\]
\end{proof}

A direct calculation yields
\[
Lh = e^{-\alpha |x|^2}\left\{4\alpha^2\sum_{i,j=1}^n a_{ij}(x)x_i x_j - 2\alpha\sum_{i=1}^n a_{ii}(x) - 2\alpha\sum_{i=1}^n b_i(x)x_i + c\right\}
\]
\[
- ce^{-\alpha r^2}
\]
\[
\ge e^{-\alpha |x|^2}\left\{4\alpha^2\sum_{i,j=1}^n a_{ij}(x)x_i x_j - 2\alpha\sum_{i=1}^n [a_{ii}(x)+b_i(x)x_i] + c\right\}.
\]
By the ellipticity, we have
\[
\sum_{i,j=1}^n a_{ij}(x)x_i x_j \ge \lambda |x|^2 \ge \lambda\left(\frac{r}{2}\right)^2 > 0 \text{ in } \Sigma.
\]
We conclude $Lh > 0$ in $\Sigma$ for $\alpha$ large enough. With such an $h$, we have $Lv = Lu + \epsilon Lh > 0$ in $\Sigma$ for any $\epsilon > 0$. By Lemma 2.1, $v$ cannot attain its nonnegative maximum inside $\Sigma$.

Next, we prove, for some small $\epsilon > 0$, $v$ attains at $x_0$ its nonnegative maximum.
Consider $v$ on the boundary $\partial\Sigma$.
\begin{enumerate}
\item[(i)] For $x \in \partial\Omega \cap B$, since $u(x) < u(x_0)$, so $u(x) < u(x_0)-\delta$ for some $\delta > 0$. Take $\epsilon$ small such that $\epsilon h < \delta$ on $\partial\Omega \cap B$. Hence, for such an $\epsilon$, we have $v(x) < u(x_0)$ for any $x \in \partial\Omega \cap B$.
\item[(ii)] On $\Sigma \cap \partial B$, $h(x)=0$ and $u(x) < u(x_0)$ for $x \neq x_0$. Hence $v(x) < u(x_0)$ on $\Sigma \cap \partial B \setminus \{x_0\}$ and $v(x_0)=u(x_0)$.
\end{enumerate}
Therefore, we conclude for any small $t > 0$
\[
\frac{1}{t}\bigl(v(x_0)-v(x_0-tv)\bigr) \ge 0.
\]
By letting $t \to 0$, we obtain
\[
\liminf_{t\to 0} \frac{1}{t}\bigl(u(x_0)-u(x_0-tv)\bigr) \ge -\epsilon \frac{\partial h}{\partial v}(x_0).
\]
By the definition of $h$, we have
\[
\frac{\partial h}{\partial v}(x_0) = \frac{\partial h}{\partial n}(x_0)\vectn\cdot v = -2\alpha r e^{-\alpha r^2}\,\vectn\cdot v < 0.
\]
This finishes the proof.

Next, we are ready to prove the strong maximum principle.

\begin{theorem}\label{hanlin:maxprin-theorem-2-8}\dcref{Theorem 2.8}\group{ch02-maximum-principles}\source{han-lin-lecture-notes:page-0031}
Let $u \in C^2(\Omega) \cap C(\overline{\Omega})$ satisfy $Lu \ge 0$ with $c(x) \le 0$ in $\Omega$.
Then the nonnegative maximum of $u$ in $\overline{\Omega}$ can be assumed only on $\partial\Omega$ unless $u$ is
a constant.
\end{theorem}
\begin{proof}
Let $M$ be the nonnegative maximum of $u$ in $\overline{\Omega}$. Set $\Sigma = \{x \in \Omega;u(x)=
M\}$. It is relatively closed in $\Omega$. We need to show $\Sigma = \Omega$.
We prove by contradiction. If $\Sigma$ is a proper subset of $\Omega$, then we may find an
open ball $B \subset \Omega \setminus \Sigma$ with a point on its boundary belonging to $\Sigma$. (In fact, we may
choose a point $p \in \Omega \setminus \Sigma$ such that $d(p,\Sigma) < d(p,\partial\Omega)$ first and then extend the ball
centered at $p$. It hits $\Sigma$ before hitting $\partial\Omega$.) Suppose $x_0 \in \partial B \cap \Sigma$. Obviously we
have $Lu \ge 0$ in $B$ and
\[
u(x) < u(x_0) \quad \text{for any } x \in B \text{ and } u(x_0)=M \ge 0.
\]
Theorem 2.6 implies $\dfrac{\partial u}{\partial n}(x_0)>0$, where $\mathbf{n}$ is the outward normal direction at $x_0$ to
the ball $B$. While $x_0$ is the interior maximal point of $\Omega$, hence $Du(x_0)=0$. This
leads to a contradiction. $\square$
\end{proof}

The following result is often referred to as the comparison principle.

\begin{corollary}\label{hanlin:maxprin-cor-2-9}\dcref{Corollary 2.9}\group{ch02-maximum-principles}\source{han-lin-lecture-notes:page-0032}
Suppose $u \in C^2(\Omega)\cap C(\overline{\Omega})$ satisfies $Lu \ge 0$ in $\Omega$ with $c(x)\le 0$
in $\Omega$. If $u \le 0$ on $\partial\Omega$, then $u \le 0$ in $\Omega$. In fact, either $u < 0$ in $\Omega$ or $u \equiv 0$ in $\Omega$.
\end{corollary}

In order to discuss the boundary value problem with a general boundary con-
dition, we need the following result, which is just a corollary of Theorem 2.6 and
Theorem 2.8.

\begin{corollary}\label{hanlin:maxprin-cor-2-10}\dcref{Corollary 2.10}\group{ch02-maximum-principles}\source{han-lin-lecture-notes:page-0032}
Suppose $\Omega$ has the interior sphere property and that $u \in$
$C^2(\Omega)\cap C^1(\overline{\Omega})$ satisfies $Lu \ge 0$ in $\Omega$ with $c(x)\le 0$. Assume u attains its nonneg-
ative maximum at $x_0 \in \overline{\Omega}$. Then $x_0 \in \partial\Omega$ and for any outward direction $v$ at $x_0$ to
$\partial\Omega$
\[
\frac{\partial u}{\partial v}(x_0)>0,
\]
\textit{unless $u$ is a constant in $\overline{\Omega}$.}
\end{corollary}

\begin{corollary}\label{hanlin:maxprin-cor-2-11}\dcref{Corollary 2.11}\group{ch02-maximum-principles}\source{han-lin-lecture-notes:page-0032}
Suppose $\Omega$ is bounded in $\mathbb{R}^n$ and satisfies the interior sphere
property. Let $u \in C^2(\Omega) \cap C^1(\overline{\Omega})$ be a solution of the following boundary value
problem
\[
Lu=f \qquad \text{in } \Omega
\]
\[
\frac{\partial u}{\partial n}+\alpha(x)u=\varphi \qquad \text{on } \partial\Omega,
\]
\textit{for some $f \in C(\overline{\Omega})$ and $\varphi \in C(\partial\Omega)$. Assume in addition that $c(x)\le 0$ in $\Omega$ and
$\alpha(x)\ge 0$ on $\partial\Omega$. Then $u$ is the unique solution if $c\not\equiv 0$ or $\alpha\not\equiv 0$. If $c\equiv 0$ and
$\alpha\equiv 0$, $u$ is unique up to additive constants.}
\end{corollary}
\begin{proof}
Suppose $u$ is a solution of the following homogeneous equation
\[
Lu=0 \qquad \text{in } \Omega
\]
\[
\frac{\partial u}{\partial n}+\alpha(x)u=0 \qquad \text{on } \partial\Omega.
\]

\textit{Case 1.} $c \not\equiv 0$ or $\alpha \not\equiv 0$. Suppose that $u$ has a positive maximum at $x_0 \in \overline{\Omega}$.
If $u$ is a positive constant, there leads to a contradiction to the condition $c\not\equiv 0$ in
$\Omega$ or $\alpha\not\equiv 0$ on $\partial\Omega$. Otherwise $x_0 \in \partial\Omega$ and $\dfrac{\partial u}{\partial n}(x_0)>0$ by Corollary 2.10, which
contradicts the boundary value. Therefore $u \equiv 0$.

\textit{Case 2.} $c \equiv 0$ and $\alpha \equiv 0$. Suppose $u$ is a nonconstant solution. Then its
maximum in $\overline{\Omega}$ is assumed only on $\partial\Omega$ by Theorem 2.8, say at $x_0 \in \partial\Omega$. Again
Corollary 2.10 implies $\dfrac{\partial u}{\partial n}(x_0)>0$. This is a contradiction. Therefore, $u$ is a
constant. $\square$
\end{proof}

The following theorem generalizes the comparison principle under no restric-
tions on $c(x)$.

\begin{theorem}\label{hanlin:maxprin-theorem-2-12}\dcref{Theorem 2.12}\group{ch02-maximum-principles}\source{han-lin-lecture-notes:page-0032,han-lin-lecture-notes:page-0033}
Suppose $u \in C^2(\Omega)\cap C(\overline{\Omega})$ satisfies $Lu \ge 0$. If $u \le 0$ in $\Omega$,
then either $u < 0$ in $\Omega$ or $u \equiv 0$ in $\Omega$.
\end{theorem}
\begin{proof}
\emph{Method 1.} Suppose $u(x_0)=0$ for some $x_0\in\Omega$. We will prove that
$u\equiv 0$ in $\Omega$. Write $c(x)=c^+(x)-c^-(x)$, where $c^+(x)$ and $c^-(x)$ are the positive
and negative part of $c(x)$ respectively. Then $u$ satisfies
\[
a_{ij}D_{ij}u+b_iD_iu-c^-u\ge -c^+u\ge 0.
\]
So we have $u=0$ by Theorem 2.8.
\emph{Method 2.} Set $v=ue^{-\alpha x_1}$ for some $\alpha>0$ to be determined. By $Lu\ge 0$, we
have
\[
a_{ij}D_{ij}v+\bigl(\alpha(a_{1i}+a_{i1})+b_i\bigr)D_iv+(a_{11}\alpha^2+b_1\alpha+c)v\ge 0.
\]
Choose $\alpha$ large enough such that $a_{11}\alpha^2+b_1\alpha+c>0$. Therefore $v$ satisfies
\[
a_{ij}D_{ij}v+\bigl(\alpha(a_{1i}+a_{i1})+b_i\bigr)D_iv\ge 0.
\]
Hence we apply Theorem 2.8 to $v$ to conclude that either $v<0$ in $\Omega$ or $v\equiv 0$ in
$\Omega$.
\end{proof}

\begin{theorem}\label{hanlin:maxprin-theorem-2-13}\dcref{Theorem 2.13}\group{ch02-maximum-principles}\source{han-lin-lecture-notes:page-0033}
Suppose there exists a $w\in C^2(\Omega)\cap C^1(\overline{\Omega})$ satisfying $w>0$ in
$\overline{\Omega}$ and $Lw\le 0$ in $\Omega$. If $u\in C^2(\Omega)\cap C(\overline{\Omega})$ satisfies $Lu\ge 0$ in $\Omega$, then
$u/w$ cannot attain in $\Omega$ its nonnegative maximum unless $u/w$ is constant. If, in addition, $u/w$
assumes its nonnegative maximum at $x_0\in\partial\Omega$ and $u/w\not\equiv \mathrm{const}.$, then for any
outward direction $\nu$ at $x_0$ to $\partial\Omega$
\[
\frac{\partial}{\partial\nu}\left(\frac{u}{w}\right)(x_0)>0,
\]
provided $\partial\Omega$ has the interior sphere property at $x_0$.
\end{theorem}
\begin{proof}
Set $v=\dfrac{u}{w}$. Then $v$ satisfies
\[
a_{ij}D_{ij}v+B_iD_iv+\left(\frac{Lw}{w}\right)v\ge 0,
\]
where $B_i=b_i+\dfrac{2}{w}a_{ij}D_{ij}w$. We simply apply Theorem 2.6 and Corollary 2.9 to
$v$.
\end{proof}

\begin{remark}\label{hanlin:maxprin-remark-2-14}\dcref{Remark 2.14}\group{ch02-maximum-principles}\source{han-lin-lecture-notes:page-0033}
If the operator $L$ in $\Omega$ satisfies the condition of Theorem 2.13,
then $L$ has the comparison principle. In particular, the Dirichlet boundary value
problem
\[
Lu=f \quad \text{in }\Omega
\]
\[
u=\varphi \quad \text{on }\partial\Omega
\]
has at most one solution.
\end{remark}

The next result is the so-called maximum principle for \emph{narrow domains}.

\begin{proposition}\label{hanlin:maxprin-prop-2-15}\dcref{Proposition 2.15}\group{ch02-maximum-principles}\source{han-lin-lecture-notes:page-0033,han-lin-lecture-notes:page-0034}
Let $d$ be a positive number and $e$ be a unit vector such that
$|(y-x)\cdot e|<d$ for any $x,y\in\Omega$. Then there exists a $d_0>0$, depending only on
$\lambda$ and the sup-norm of $b_i$ and $c^+$, such that the assumptions of Theorem 2.13 are
satisfied if $d\le d_0$.
\end{proposition}
\begin{proof}
By choosing $e=(1,0,\cdots,0)$, we suppose $\bar{\Omega}$ lies in $\{0<x_1<d\}$. Assume, in addition,
$|b_i|,c^+\le N$ for some positive constant $N$. We construct $w$ as
follows. Set $w=e^{\alpha d}-e^{\alpha x_1}>0$ in $\bar{\Omega}$. A direct calculation yields
\[
Lw=-(a_{11}\alpha^2+b_1\alpha)e^{\alpha x_1}+c(e^{\alpha d}-e^{\alpha x_1})
\le -(a_{11}\alpha^2+b_1\alpha)+Ne^{\alpha d}.
\]
Choose $\alpha$ large to have
\[
a_{11}\alpha^2+b_1\alpha\ge \lambda\alpha^2-N\alpha\ge 2N.
\]
Hence $Lw\le -2N+Ne^{\alpha d}=N(e^{\alpha d}-2)\le 0$ if $d$ is small with $e^{\alpha d}\le 2$.
\hfill$\square$
\end{proof}

\begin{remark}\label{hanlin:maxprin-remark-2-16}\dcref{Remark 2.16}\group{ch02-maximum-principles}\source{han-lin-lecture-notes:page-0034}
Some results in this section, including Proposition 2.15, hold
for unbounded domains.
\end{remark}

\subsection*{2.2. Poisson Equations}

In this section, we use the maximum principle to discuss solutions of Poisson
equations $\Delta u=f$.

As the first application, we derive interior gradient estimates for solutions of
Poisson equations.

\begin{lemma}\label{hanlin:poisson-lemma-2-17}\dcref{Lemma 2.17}\group{ch02-maximum-principles}\source{han-lin-lecture-notes:page-0034,han-lin-lecture-notes:page-0035}
Suppose $u\in C^2(B_R)\cap C(\bar{B}_R)$ satisfies
\[
\Delta u=f \text{ in } B_R,
\]
for some $f\in C(\bar{B}_R)$. Then
\[
|Du(0)|\le \frac{n}{R}\max_{\partial B_R}|u|+\frac{R}{2}\max_{B_R}|f|.
\]
\end{lemma}
\begin{proof}
We write $B=B_R$ and consider $D_nu(0)$. Set $M=\max_{\partial B}|u|$ and
$F=\max_B|f|$. Consider
\[
v(x',x_n)=\frac12\bigl(u(x',x_n)-u(x',-x_n)\bigr)\quad\text{in }B^+.
\]
Then $v$ satisfies
\[
|\Delta v|\le F \quad\text{in }B^+,
\]
\[
|v|\le M \quad\text{on }\partial B^+,
\]
\[
v(x',0)=0\quad\text{for }|x'|<R.
\]
Consider an auxiliary function
\[
w(x',x_n)=A|x'|^2+Bx_n+Cx_n^2.
\]
By choosing $A,B,C$ appropriately, we assume
\begin{align*}
(\text{i})\;&w(x',0)=A|x'|^2\ge 0 \quad\text{for }|x'|\le R;\\
(\text{ii})\;&w(x',x_n)=AR^2+(C-A)x_n^2+Bx_n\ge M \quad\text{for }|x'|^2+x_n^2=R^2;\\
(\text{iii})\;&\Delta w=2(n-1)A+2C\le -F \quad\text{in }B^+.
\end{align*}
To achieve this, we need to require
\[
A\ge 0,
\]
\[
AR^2\ge M,
\]
\[
(C-A)x_n^2+Bx_n\ge 0 \quad\text{for any }0\le x_n\le R,
\]
\[
2(n-1)A+2C\le -F.
\]
So we take
\[
A = \frac{M}{R^2},
\]
and
\[
C = -\frac{F}{2} - (n-1)A,
\]
which implies
\[
A - C = \frac{F}{2} + nA.
\]
For $B$, we need $B + (C-A)x_n \ge 0$, or
\[
B \ge \left(\frac{F}{2} + \frac{n}{R^2}M\right)x_n \qquad \text{for any $0 \le x_n \le R$.}
\]
Hence we may take
\[
B = \left(\frac{F}{2} + \frac{n}{R^2}M\right)R = \frac{R}{2}F + \frac{n}{R}M.
\]
With such a choice of $w$, we have $\Delta w \le \Delta v$ in $B^+$ and $w \ge v$ on $\partial B^+$. By Corollary 2.9, the comparison principle, we obtain $w \ge v$ in $B^+$. Taking $x' = 0$, we obtain
\[
\frac{1}{2}\left(\frac{u(0,x_n)-u(0,-x_n)}{x_n}\right) \le B + Cx_n \qquad \text{for any $0 < x_n < R$.}
\]
Therefore, we have by letting $x_n \to 0$
\[
D_n u(0) \le \frac{n}{R}M + \frac{R}{2}F.
\]
By considering $-u$ similarly, we obtain
\[
|D_n u(0)| \le \frac{n}{R}M + \frac{R}{2}F.
\]
This finishes the proof.
\end{proof}

\begin{lemma}\label{hanlin:poisson-lemma-2-18}\dcref{Lemma 2.18}\group{ch02-maximum-principles}\source{han-lin-lecture-notes:page-0035,han-lin-lecture-notes:page-0036}
Suppose $u \in C^2(B_{2R}) \cap C(\bar B_{2R})$ satisfies
\[
\Delta u = f \quad \text{in } B_{2R},
\]
for some $f \in C(\bar B_{2R})$. Then for any $x,y \in B_R$ with $x \neq y$,
\[
R^2 \left|\frac{Du(x)-Du(y)}{|x-y|}\right| \le c \left(\sup_{B_{2R}} |u| + R^2 \sup_{B_{2R}} |f|\right)\left(\log \frac{2R}{|x-y|} + 1\right),
\]
where $c$ is a positive constant depending only on $n$.
\end{lemma}
\begin{proof}
In the following, we use cubes $Q_R, Q_{2R}$ instead of balls $B_R, B_{2R}$. Set
\[
M = \sup_{Q_{2R}} |u| \qquad \text{and} \qquad F = \sup_{Q_{2R}} |f|.
\]
We will prove for any $x_n \in \left(0,\frac{R}{4}\right)$,
\begin{equation}
\frac{1}{2}|Du(0,x_n)-Du(0,-x_n)| \le c x_n \left(\frac{M}{R^2} + F\right)\left(\log \frac{2R}{x_n} + 1\right),
\tag{1}
\end{equation}
where $c$ is a positive constant depending only on $n$. By Lemma 2.17, we have
\begin{equation}
\sup_{Q_R} |Du| \le c\left(\frac{M}{R} + RF\right).
\tag{2}
\end{equation}
\end{proof}
% Source page 0036
Let $Q'$ be the domain in $\R^{n+1}$ given by
\[
Q'=\left\{(x_1,\ldots,x_{n-1},y,z);\ |x_i|<\frac{R}{2}, i=1,\ldots,n-1,\ 0<y,z<\frac{R}{4}\right\},
\]
and define in $Q'$ the function
\[
v(x',y,z)=\frac{1}{4}\{u(x',y+z)-u(x',y-z)-u(x',-y+z)+u(x',-y-z)\}.
\]
Define an operator $L$ in $\R^{n+1}$ by
\[
L\equiv \sum_{i=1}^{n-1} D_{x_i x_i}+\frac{1}{2}D_{yy}+\frac{1}{2}D_{zz}.
\]

It is easy to see in $Q'$
\begin{align*}
(\mathrm{i})\quad &|Lu|\leq F \text{ in } Q';\\
(\mathrm{ii})\quad &v(x',0,z)=v(x',y,0)=0;\\
(\mathrm{iii})\quad &|v|\leq M \text{ on } |x_i|=\frac{R}{2}, i=1,\ldots,n-1;\\
(\mathrm{iv})\quad &v(x',\tfrac{R}{4},z)\leq \mu z \text{ and } v(x',y,\tfrac{R}{4})\leq \mu y,
\end{align*}
where $|Du|\leq \mu$ in $Q_R$ with $\mu$ given in terms of $M$ and $F$ by (2). Choose a comparison function $w$ in $Q'$ of the form
\[
w(x',y,z)=\frac{4M|x'|^2}{R^2}+\frac{4\mu}{R}yz+kyz\log\frac{2R}{y+z},
\]
where $k$ is a positive constant to be determined. Note first that $|v|\leq w$ on $\partial Q'$.
Since
\[
Lw(x',y,z)=\frac{8(n-1)}{R^2}M+k\left(-1+\frac{yz}{(y+z)^2}\right)\leq \frac{8(n-1)}{R^2}M-\frac{3}{4}k,
\]
we see that $Lw\leq -F$ provided
\[
k\geq \frac{4}{3}\left(F+\frac{8(n-1)M}{R^2}\right).
\]

With such a choice of $k$, the function
\[
w(x',y,z)=\frac{4M|x'|^2}{R^2}+yz\left(\frac{4\mu}{R}+k\log\frac{2R}{y+z}\right)
\]
satisfies the conditions $L(w\pm v)\leq 0$ in $Q'$ and $w\pm v\geq 0$ on $\partial Q'$. By Corollary 2.9, the comparison principle, $|v|\leq w$ in $Q'$. Letting $x'=0$ in this inequality, then dividing by $z$ and letting $z$ tend to zero, we obtain
\[
\frac{1}{2}|D_yu(0,y)-D_yu(0,-y)|\leq \frac{4\mu}{R}y+ky\log\frac{2R}{y}.
\]

We proved (1) for $D_n$.

With a slight modification we can derive (1) for $D_i, i=1,\ldots,n-1$. We work in $\R^n$ in this case. In the domain
\[
Q'=\left\{(x_1,\ldots,x_{n-2},y,z);\ |x_i|<\frac{R}{2}, i=1,\ldots,n-2, 0<y,z<\frac{R}{2}\right\},
\]
we define
\[
v(\hat{x},y,z)=\frac{1}{4}\{u(\hat{x},y,z)-u(\hat{x},-y,-z)-u(\hat{x},y,-z)+u(\hat{x},-y,-z)\}.
\]
% Source page 0037
We choose the comparison function of the form
\[
w(\hat{x}, y, z)
= \frac{4M|\hat{x}|^2}{R^2}
  + yz\left(\frac{4\mu}{R} + \bar{k}\log\frac{2R}{y+z}\right),
\]
with $\mu$ as before and
\[
\bar{k} \ge \frac{2}{3}\left(F + \frac{8(n-2)M}{R^2}\right).
\]

We verify easily that $\Delta(w \pm v) \le 0$ in $Q'$ and $w \pm v \ge 0$ on $\partial Q'$. Hence $|v| \le w$ in $Q'$. As above, if we set $\hat{x} = 0$, then divide by $y$ and let $y$ tend to zero, we obtain
\[
\frac{1}{2}\bigl|D_{n-1}u(0,z) - D_{n-1}u(0,-z)\bigr|
\le \frac{4\mu}{R} z + \bar{k}z\log\frac{2R}{z}.
\]
Obviously the same result holds if $D_{n-1}$ is replaced by $D_i$ for $i = 1,\ldots,n-2$. $\square$

Despite the elementary character of its proof, Lemma 2.18 is essentially sharp and the estimate cannot be improved without further continuity assumptions on $f$.

As the second application, we discuss the Schauder theory. We will show that, if $u \in C^2(B_1)$ satisfies
\[
\Delta u = f \quad \text{in } B_1,
\]
for some Hölder continuous function $f$ in $B_1$, then $D^2u$ is Hölder continuous with the same exponent. We will use the maximum principle approach, avoiding the potential integrals.

Let us recall the definition of the Hölder continuity. A function $u$ is $C^\alpha$ at $0$ if
\[
|u(x) - u(0)| \le c|x|^\alpha.
\]
Here $|x|^\alpha$ is of course smaller than any constants. It gives a quantitative speed of how functions approach constant. We define the Hölder semi-norm of $u$ at $0$ by
\[
[u]_{C^\alpha}(0) \equiv \sup_{|x|\le 1}\frac{|u(x) - u(0)|}{|x|^\alpha}.
\]

Similarly, we can define $C^{1,\alpha}$ and $C^{2,\alpha}$. For example, $u$ is $C^{2,\alpha}$ at $0$ if there exists a second order polynomial $P(x)$ such that
\[
|u(x) - P(x)| \le c|x|^{2+\alpha}.
\]

\begin{lemma}\label{hanlin:maxprin-lemma-2-19}\dcref{Lemma 2.19}\group{ch02-maximum-principles}\source{han-lin-lecture-notes:page-0037,han-lin-lecture-notes:page-0038}
Suppose $u \in C^2(B_1)$ satisfies
\[
\Delta u = f \quad \text{in } B_1,
\]
for some $f \in C(B_1)$. Then for any $\alpha \in (0,1)$, there exist constants $c_0 > 0$, $\mu \in (0,1)$ and $\epsilon_0 > 0$, depending only on $n$ and $\alpha$, such that, if $|u| \le 1$ and $|f| \le \epsilon_0$ in $B_1$, there exists a second order harmonic polynomial
\[
p(x) = \frac{1}{2}x^T A x + B \cdot x + C,
\]
satisfying
\[
|u(x) - p(x)| \le \mu^{2+\alpha} \quad \text{for } |x| \le \mu,
\]
and
\[
|A| + |B| + |C| \le c_0.
\]
\end{lemma}
\begin{proof}
Suppose $v$ is the harmonic function satisfying
\[
\Delta v = 0 \text{ in } B_1,
\qquad
v = u \text{ on } \partial B_1.
\]
Then the function $u-v$ satisfies
\[
\Delta (u-v) = f \text{ in } B_1,
\qquad
u-v = 0 \text{ on } \partial B_1.
\]
By Corollary 2.9, the comparison principle, we have
\[
|u(x)-v(x)| \le \frac{1-|x|^2}{2n}\sup_{B_1}|f| \qquad \text{for any } x \in B_1.
\]
Clearly $|v| \le 1$ in $B_1$. Hence its second order Taylor polynomial at $0$
\[
p(x)=\frac12 x^T A x + B \cdot x + C
\]
has universal bounded coefficients and is harmonic. By the mean value theorem,
we have
\[
|u(x)-p(x)| \le |v(x)-p(x)|+\frac1{2n}\sup_{B_1}|f|
\]
\[
\le C|x|^3\sup_{B_{1/2}}|D^3v|+\frac1{2n}\sup_{B_1}|f|
\qquad \text{for any } x \in B_{1/2}.
\]
Now take $\mu$ small enough such that the first term is less than or equal to
\[
\frac12\mu^{2+\alpha} \qquad \text{for } |x|\le \mu,
\]
and then take $\varepsilon_0$ such that
\[
\frac12\sup_{B_1}|f| \le \frac12\varepsilon_0 \le \frac12\mu^{2+\alpha}.
\]
This finishes the proof.
\end{proof}

\begin{theorem}\label{hanlin:maxprin-theorem-2-20}\dcref{Theorem 2.20}\group{ch02-maximum-principles}\source{han-lin-lecture-notes:page-0038,han-lin-lecture-notes:page-0039}
Suppose $u \in C^2(B_1)$ satisfies
\[
\Delta u = f \text{ in } B_1,
\]
where $f$ is H\"older continuous at $0$. Then $u(x)$ is $C^{2,\alpha}$ at $0$, i.e., there exists a second order polynomial
\[
p(x)=\frac12 x^T A x + Bx + C
\]
satisfying
\[
|u(x)-p(x)| \le c_0|x|^{2+\alpha}\bigl(|u|_{L^\infty} + |f(0)| + [f]_{C^\alpha(0)}\bigr) \qquad \text{for any } x \in B_1,
\]
and
\[
|A| + |B| + |C| \le c_0\bigl(|u|_{L^\infty} + |f(0)| + [f]_{C^\alpha(0)}\bigr),
\]
where $c_0$ is a positive constant depending only on $n$ and $\alpha$.
\end{theorem}
% Source page 0038
\begin{proof}\uses{hanlin:maxprin-lemma-2-19}
Without loss of generality, we assume $f(0)=0$. In general, we set
\[
v=u-\frac{1}{2n}f(0)|x|^2.
\]
Then $\Delta v=f-f(0)$ in $B_1$. Furthermore, we assume $|u|\le 1$ and
$[f]_{C^\alpha}(0)\le \epsilon_0$ for small $\epsilon_0>0$. The general case can be recovered by considering
\[
\frac{u}{|u|_{L^\infty}+\frac{1}{\epsilon_0}[f]_{C^\alpha}(0)}.
\]

First we claim that there are harmonic polynomials for any $k=1,2,\cdots,$
\[
P_k(x)=\frac{1}{2}x^T A_k x+B_k x+C_k,
\]
satisfying
\[
|u(x)-P_k(x)|\le \mu^{(2+\alpha)k}\qquad \text{for any } |x|\le \mu^k,
\]
and
\[
|A_k-A_{k+1}|\le c\mu^{\alpha k},
\]
\[
|B_k-B_{k+1}|\le c\mu^{(\alpha+1)k},
\]
\[
|C_k-C_{k+1}|\le c\mu^{(\alpha+2)k},
\]
where $\mu\in(0,1)$ and $c$ are positive constants depending only on $n$ and $\alpha$.
Note that the case $k=1$ corresponds to Lemma 2.19. Let us assume it is true
for $k$. Set
\[
w(y)=\frac{(u-P_k)(\mu^k y)}{\mu^{(2+\alpha)k}}\qquad \text{for } |y|\le 1.
\]
Then we have
\[
\Delta w(y)=\frac{f(\mu^k y)}{\mu^{\alpha k}}\qquad \text{for } y\in B_1.
\]
By Lemma 2.19, there is a harmonic polynomial $p_0$ with bounded coefficients such
that
\[
|w(y)-p_0(y)|\le \mu^{2+\alpha}\qquad \text{for } |y|\le \mu,
\]
provided
\[
\sup_{|y|\le 1}\frac{|f(\mu^k y)|}{\mu^{\alpha k}}\le [f]_{C^0}(0)\le \epsilon_0.
\]

Now we scale back to get
\[
|u(x)-P_k(x)-\mu^{(2+\alpha)k}p_0\!\left(\frac{x}{\mu^k}\right)|\le \mu^{(k+1)(2+\alpha)}\qquad \text{for } |x|\le \mu^{k+1}.
\]
Clearly we proved the $(k+1)$-th step by letting
\[
P_{k+1}(x)=P_k(x)+\mu^{(2+\alpha)k}p_0\!\left(\frac{x}{\mu^k}\right).
\]
It is easy to see that $A_k$, $B_k$ and $C_k$ converge and the limiting polynomial
\[
p(x)=\frac{1}{2}x^T A_\infty x+B_\infty x+C_\infty
\]
satisfies
\[
|P_k(x)-p(x)|\le c\{ |x|^2\mu^{\alpha k}+|x|\mu^{(\alpha+1)k}+\mu^{(\alpha+2)k}\}\le c\mu^{(2+\alpha)k},
\]
for any $|x|\le \mu^k$. Hence we have for $|x|\le \mu^k$
\[
|u(x)-p(x)|\le |u(x)-P_k(x)|+|P_k(x)-p(x)|\le c\mu^{(\alpha+2)k},
\]
% Source page 0039
or
\[
\lvert u(x)-p(x)\rvert\le c\lvert x\rvert^{2+\alpha}\qquad \text{for any }x\in B_1.
\]
This finishes the proof.
\end{proof}

As the third application, we use the moving plane method to discuss the symmetry of solutions. The following result was first proved by Gidas, Ni and Nirenberg.

\begin{theorem}\label{hanlin:maxprin-theorem-2-21}\dcref{Theorem 2.21}\group{ch02-maximum-principles}\source{han-lin-lecture-notes:page-0040,han-lin-lecture-notes:page-0041}
\emph{Suppose $u\in C^2(\overline{B}_1)$ is a positive solution of}
\[
\Delta u=f(u)\ \text{ in }B_1,
\]
\[
u=0\ \text{ on }\partial B_1,
\]
\emph{where $f$ is locally Lipschitz in $\mathbb{R}$. Then $u$ is radially symmetric in $B_1$ and $\dfrac{\partial u}{\partial r}(x)<0$}
\emph{for $x\neq 0$.}
\end{theorem}

We need a lemma first.

\begin{lemma}\label{hanlin:maxprin-lemma-2-22}\dcref{Lemma 2.22}\group{ch02-maximum-principles}\source{han-lin-lecture-notes:page-0040}\uses{hanlin:maxprin-theorem-2-21}
\emph{Let $u$ be as in Theorem 2.21. Then for any $x_0\in\partial B_1$ and a unit vector $\nu$ with $\nu\cdot x_0>0$, there exists an $r>0$ such that}
\[
\frac{\partial u}{\partial \nu}(x)<0\qquad \text{for any }x\in B_1\cap B_r(x_0).
\]
\end{lemma}

\begin{proof}\uses{hanlin:maxprin-theorem-2-6}
We consider two cases.
(i) $f(0)\le 0$. Then we have
\[
\Delta u=f(u)\le f(u)-f(0)\le \left(\frac{f(u)-f(0)}{u}\right)u(x)=c(x)u(x),
\]
where $c(x)$ is a nonnegative bounded function. Theorem 2.6, the Hopf Lemma, implies $\dfrac{\partial u}{\partial \nu}(x_0)<0$.

(ii) $f(0)>0$. Clearly, we have $\dfrac{\partial u}{\partial r}(x)\le 0$ for any $x\in\partial B_1$. Suppose there exists a sequence $x_i\to x_0\in\partial B_1$ such that
\[
\frac{\partial u}{\partial r}(x_i)\ge 0.
\]
Then $\dfrac{\partial u}{\partial r}(x_0)=0$. Hence $\dfrac{\partial u}{\partial r}$ is maximized at $x_0$ along $\partial B_1$, and then
\[
\frac{\partial^2u}{\partial T\partial r}(x_0)=0,
\]
for any tangential direction $T$. Clearly $\dfrac{\partial^2u}{\partial T^2}(x_0)=0$. Therefore we have
\[
\frac{\partial^2u}{\partial n^2}(x_0)=\Delta u(x_0)=f(0)>0.
\]

We may assume $x_0=(1,0,\cdots,0)$ and then apply Taylor expansion to $\dfrac{\partial u}{\partial n}(x)$ at $x_0$. Then we get
\[
\frac{\partial u}{\partial n}(x)=\frac{\partial u}{\partial n}(x_0)+D\frac{\partial u}{\partial n}(x_0)\cdot(x-x_0)+o(\lvert x-x_0\rvert)
\]
\[
=\frac{\partial^2u}{\partial n^2}(x_0)(x_1-1)+o(\lvert x-x_0\rvert).
\]
% Source page 0040
Hence for $x_1<1$ small, we have $\dfrac{\partial u}{\partial n}(x)<0$. This contradiction finishes the proof. \(\square\)
\end{proof}

\begin{proof}\uses{hanlin:maxprin-lemma-2-22,hanlin:maxprin-theorem-2-13,hanlin:maxprin-theorem-2-6}
We only need to show that $u$ is symmetric in $x_1$ direction. Define for $\lambda>0$
\[
\Sigma_\lambda=\{x\in B_1;x_1>\lambda\},
\qquad
T_\lambda=\{x_1=\lambda\},
\]
\[
\Sigma_\lambda'=\text{the reflection of }\Sigma_\lambda\text{ with respect to }T_\lambda,
\qquad
x_\lambda=(2\lambda-x_1,x_2,\cdots,x_n).
\]
First we write the following statement
\[
(*)
\qquad
u(x)<u(x_\lambda)\ \text{ for any }x\in\Sigma_\lambda\text{ and }u_{x_1}<0\text{ on }B_1\cap T_\lambda.
\]
We claim that $\Lambda=\{\lambda\in(0,1);(*)\text{ holds for }\lambda\}$ is $(0,1)$.

To this end, we need to show that $\Lambda$ is non-empty, open and closed in $(0,1)$. It is clear that $\Lambda$ is non-empty by Lemma 2.22. For the openness, we take a $\lambda_0\in\Lambda$ and prove that $(*)$ holds for $\lambda<\lambda_0$ if $\lambda$ is closed to $\lambda_0$. First we have
\[
u(x)<u(x_{\lambda_0})\qquad \text{for any }x\in\Sigma_{\lambda_0}.
\]
Take any $\lambda_1>\lambda_0$ with $\lambda_1$ closed to $\lambda_0$. Then we get
\[
u(x)<u(x_{\lambda_0})\qquad \text{for any }x\in\overline{\Sigma}_{\lambda_1}.
\]
By the continuity, this implies for $\lambda$ close to $\lambda_0$
\[
u(x)<u(x_\lambda)\qquad \text{for any }x\in\overline{\Sigma}_{\lambda_1}.
\]
Next, Lemma 2.22 and the fact
\[
D_{x_1}u(x)<0\ \text{ on }B_1\cap T_{\lambda_0}
\]
imply that
\[
D_{x_1}u(x)<0\ \text{ on }B_1\cap T_\lambda,
\]
for any $\lambda$ close to $\lambda_0$. Therefore, we have $\lambda\in\Lambda$ for any $\lambda$ close to $\lambda_0$. For the closedness, we prove that positive $\lambda\in\Lambda$ if $u(x)\leq u(x_\lambda)$ for $x\in\Sigma_\lambda$ and $D_{x_1}u(x)\leq 0$ for $x\in T_\lambda$. For such a $\lambda$, define $v(x)=u(x_\lambda)$ in $\Sigma_\lambda'$. Then, we have
\[
\Delta v=f(v)\ \text{ in }\Sigma_\lambda'.
\]
Consider $w=u-v$ in $\Sigma_\lambda'$. We have $w\geq 0$ and $w\neq 0$ in $\Sigma_\lambda'$. In fact, $w>0$ on $\partial\Sigma_\lambda'\setminus T_\lambda$. Moreover, we get
\[
\Delta w=f(u)-f(v)=c(x)w,
\]
by the mean value theorem. Theorem 2.13 implies $w>0$ in $\Sigma_\lambda$. We note $w=0$ on $T_\lambda$. By Theorem 2.6, the Hopf lemma, we get
\[
2D_{x_1}u=D_{x_1}w<0\ \text{ on }T_\lambda\cap B_1.
\]
Hence $\lambda\in\Lambda$. This finishes the proof of Theorem 2.21. \(\square\)
\end{proof}

\begin{remark}\label{hanlin:maxprin-remark-2-23}\dcref{Remark 2.23}\group{ch02-maximum-principles}\source{han-lin-lecture-notes:page-0041}
The method above depends on the smoothness of domains and the smoothness of solutions up to the boundary. In fact, such conditions can be removed. See Section 2.6 for details.
\end{remark}

% Source page 0041
\section*{2.3. A Priori Estimates}

In this section, we derive a priori estimates for solutions of the Dirichlet problem
and the Neumann problem.
Suppose $\Omega$ is a bounded and connected domain in $\mathbb{R}^n$. Consider the operator
$L$ in $\Omega$
\[
Lu \equiv a_{ij}(x)D_{ij}u + b_i(x)D_iu + c(x)u,
\]
for $u \in C^2(\Omega)\cap C(\overline{\Omega})$. We assume that $a_{ij}$, $b_i$ and $c$ are continuous and hence
bounded in $\overline{\Omega}$ and that $L$ is uniformly elliptic in $\Omega$, i.e.,
\[
a_{ij}(x)\xi_i\xi_j \ge \lambda |\xi|^2 \qquad \text{for any } x\in \Omega \text{ and any } \xi \in \mathbb{R}^n,
\]
where $\lambda$ is a positive constant. We denote by $\Lambda$ the sup-norms of $a_{ij}$ and $b_i$, i.e.,
\[
\max_{\Omega}|a_{ij}| + \max_{\Omega}|b_i| \le \Lambda.
\]

\begin{theorem}\label{hanlin:maxprin-theorem-2-24}\dcref{Theorem 2.24}\group{ch02-maximum-principles}\source{han-lin-lecture-notes:page-0042}
Suppose $u \in C^2(\Omega)\cap C(\overline{\Omega})$ satisfies
\[
Lu = f \quad \text{in } \Omega,
\]
\[
u = \varphi \quad \text{on } \partial\Omega,
\]
for some $f \in C(\overline{\Omega})$ and $\varphi \in C(\partial\Omega)$. If $c(x)\le 0$, then
\[
\max_{\Omega}|u| \le \max_{\partial\Omega}|\varphi| + C\max_{\Omega}|f|,
\]
where $C$ is a positive constant depending only on $\lambda$, $\Lambda$ and $\operatorname{diam}(\Omega)$.
\end{theorem}

\begin{proof}\uses{hanlin:maxprin-cor-2-9}
We will construct a function $w$ in $\Omega$ such that
\begin{enumerate}
\item[(i)] $L(w\pm u)=Lw\pm f\le 0$, or $Lw\le \mp f$ in $\Omega$;
\item[(ii)] $w\pm u=w\pm \varphi \ge 0$, or $w\ge \mp\varphi$ on $\partial\Omega$.
\end{enumerate}
Set $F=\max_{\Omega}|f|$ and $\Phi=\max_{\partial\Omega}|\varphi|$. We need
\[
Lw \le -F \qquad \text{in } \Omega
\]
\[
w \ge \Phi \qquad \text{on } \partial\Omega.
\]

Suppose the domain $\Omega$ lies in the set $\{0<x_1<d\}$ for some $d>0$. Set for some
$\alpha>0$ to be chosen later
\[
w=\Phi + (e^{\alpha d}-e^{\alpha x_1})F.
\]
Then we have by a direct calculation
\[
-Lw=(a_{11}\alpha^2+b_1\alpha)Fe^{\alpha x_1}-c\Phi-c(e^{\alpha d}-e^{\alpha x_1})F
\]
\[
\ge (a_{11}\alpha^2+b_1\alpha)F \ge (\alpha^2\lambda+b_1\alpha)F \ge F,
\]
by choosing $\alpha$ large such that $\alpha^2\lambda+b_1(x)\alpha\ge 1$ for any $x\in\Omega$. Hence $w$ satisfies
(i) and (ii). By Corollary 2.9, the comparison principle, we conclude $-w\le u\le w$
in $\Omega$, and in particular,
\[
\sup_{\Omega}|u| \le \Phi + (e^{\alpha d}-1)F,
\]
where $\alpha$ is a positive constant depending only on $\lambda$ and $\Lambda$.
\hfill$\square$
\end{proof}

% Source page 0042
\begin{theorem}\label{hanlin:maxprin-theorem-2-25}\dcref{Theorem 2.25}\group{ch02-maximum-principles}\source{han-lin-lecture-notes:page-0043,han-lin-lecture-notes:page-0044}
Suppose $u \in C^2(\Omega) \cap C^1(\bar{\Omega})$ satisfies
$$
Lu = f \quad \text{in } \Omega,
$$
$$
\frac{\partial u}{\partial n} + \alpha(x)u = \varphi \quad \text{on } \partial \Omega,
$$
where $n$ is the outward normal direction to $\partial \Omega$. If $c(x) \leq 0$ in $\Omega$ and $\alpha(x) \geq \alpha_0 > 0$ on $\partial \Omega$, then
$$
\max_{\partial \Omega} |u| \leq C \bigl( \max_{\Omega} |\varphi| + \max_{\Omega} |f| \bigr),
$$
where $C$ is a positive constant depending only on $\lambda$, $\Lambda$, $\alpha_0$ and $\operatorname{diam}(\Omega)$.
\end{theorem}

\begin{proof}
We first consider a special case $c(x) \leq -c_0 < 0$. We will show
$$
|u(x)| \leq \frac{1}{c_0} F + \frac{1}{\alpha_0} \Phi \qquad \text{for any } x \in \Omega.
$$
Set
$$
v = \frac{1}{c_0} F + \frac{1}{\alpha_0} \Phi \pm u.
$$
Then we have
$$
Lv = c(x)\left(\frac{1}{c_0}F + \frac{1}{\alpha_0}\Phi\right) \pm f \leq -F \pm f \leq 0 \quad \text{in } \Omega,
$$
and
$$
\frac{\partial v}{\partial n} + \alpha v = \alpha \left(\frac{1}{c_0}F + \frac{1}{\alpha_0}\Phi\right) \pm \varphi \geq \Phi \pm \varphi \geq 0 \quad \text{on } \partial \Omega.
$$
If $v$ has a negative minimum in $\bar{\Omega}$, then $v$ attains it on $\partial \Omega$, say at $x_0 \in \partial \Omega$, by Theorem 2.3. This implies $\frac{\partial v}{\partial n}(x_0) \leq 0$ for $n = n(x_0)$, the outward normal direction at $x_0$. Therefore, we get
$$
\left(\frac{\partial v}{\partial n} + \alpha v\right)(x_0) \leq \alpha v(x_0) < 0,
$$
which is a contradiction. Hence we have $v \geq 0$ in $\bar{\Omega}$, and in particular,
$$
|u(x)| \leq \frac{1}{c_0} F + \frac{1}{\alpha_0} \Phi \qquad \text{for any } x \in \Omega.
$$
Note that $c_0$ and $\alpha_0$ are independent of $\lambda$ and $\Lambda$ for this special case.

Next, we consider the general case $c(x) \leq 0$ for any $x \in \Omega$. Consider an auxiliary function $u(x) = z(x)w(x)$ where $z$ is a positive function in $\bar{\Omega}$ to be determined. A direct calculation shows that $w$ satisfies
$$
a_{ij}D_{ij}w + B_iD_iw + \left(c + \frac{a_{ij}D_{ij}z + b_iD_iz}{z}\right) w = \frac{f}{z} \quad \text{in } \Omega,
$$
$$
\frac{\partial w}{\partial n} + \left(\alpha + \frac{1}{z}\frac{\partial z}{\partial n}\right) w = \frac{\varphi}{z} \quad \text{on } \partial \Omega,
$$
where
$$
B_i = \frac{1}{z}(a_{ij} + a_{ji})D_jz + b_i.
$$
% Source page 0043
% Source page 0044
We need to choose the function $z>0$ in $\bar\Omega$ such that
\[
c+\frac{a_{ij}D_{ij}z+b_iD_iz}{z}\le -c_0 \quad \text{in } \Omega,
\]
\[
\alpha+\frac{1}{z}\frac{\partial z}{\partial n}\ge \frac{1}{2}\alpha_0 \quad \text{on } \partial\Omega,
\]
or
\[
\frac{a_{ij}D_{ij}z+b_iD_iz}{z}\le -c_0 \quad \text{in } \Omega,
\]
\[
\left|\frac{1}{z}\frac{\partial z}{\partial n}\right|\le \frac{1}{2}\alpha_0 \quad \text{on } \partial\Omega.
\]
where $c_0$ is a positive constant depending only on $\lambda,\Lambda,\alpha_0$ and $\operatorname{diam}\Omega$. Suppose the domain $\Omega$ lies in $\{0<x_1<d\}$. Choose $z(x)=A+e^{\beta d}-e^{\beta x_1}$ for $x\in\Omega$, for some positive $A$ and $\beta$ to be determined. A direct calculation shows
\[
-\frac{1}{z}\left(a_{ij}D_{ij}z+b_iD_iz\right)
=\frac{(\beta^2a_{11}+\beta b_1)e^{\beta x_1}}{A+e^{\beta d}-e^{\beta x_1}}
\]
\[
\ge \frac{\beta^2a_{11}+\beta b_1}{A+e^{\beta d}}\ge \frac{1}{A+e^{\beta d}}>0,
\]
if $\beta$ is chosen such that $\beta^2a_{11}+\beta b_1\ge 1$. Then we have
\[
\left|\frac{1}{z}\frac{\partial z}{\partial n}\right|\le \frac{\beta}{A}e^{\beta d}\le \frac{1}{2}\alpha_0,
\]
if $A$ is chosen large. This reduces to the special case we just discussed. The new extra first order term does not change the result. We may apply the special case to $w$.

\end{proof}

\begin{remark}\label{hanlin:maxprin-remark-2-26}\dcref{Remark 2.26}\group{ch02-maximum-principles}\source{han-lin-lecture-notes:page-0044}
The result fails if we just assume $\alpha(x)\ge 0$ on $\partial\Omega$. In fact, we cannot even get the uniqueness.
\end{remark}

\section*{2.4. Gradient Estimates}

The basic method to derive gradient estimates, the so-called the Bernstein method, involves a differentiation of the equation with respect to $x_k$, $k=1,\ldots,n$, followed by a multiplication by $D_k u$ and summation over $k$. The maximum principle is then applied to the resulting equation in the function $v=|Du|^2$, possibly with some modifications. There are two kinds of gradient estimates, global gradient estimates and interior gradient estimates. We will use semi-linear equations to illustrate the idea.

Suppose $\Omega$ is a bounded and connected domain in $\mathbb R^n$. Consider the equation
\[
a_{ij}(x)D_{ij}u+b_i(x)D_i u=f(x,u) \quad \text{in } \Omega,
\]
for $u\in C^2(\Omega)$ and $f\in C(\Omega\times\mathbb R)$. We always assume that $a_{ij}$ and $b_i$ are continuous and hence bounded in $\bar\Omega$ and that the equation is uniformly elliptic in $\Omega$ in the following sense
\[
a_{ij}(x)\xi_i\xi_j\ge \lambda|\xi|^2 \quad \text{for any } x\in \Omega \text{ and any } \xi\in \mathbb R^n,
\]
for some positive constant $\lambda$.
We first derive a global gradient estimate.

\begin{theorem}\label{hanlin:maxprin-theorem-2-27}\dcref{Theorem 2.27}\group{ch02-maximum-principles}\source{han-lin-lecture-notes:page-0044,han-lin-lecture-notes:page-0045}
Suppose $u \in C^3(\Omega) \cap C^1(\bar{\Omega})$ satisfies
\begin{equation}
a_{ij}(x)D_{ij}u+b_i(x)D_i u=f(x,u) \quad \text{in } \Omega,
\tag{1}
\end{equation}
for $a_{ij}, b_i \in C^1(\bar{\Omega})$ and $f \in C^1(\bar{\Omega}\times \mathbb R)$. Then
\[
\sup_{\Omega} |Du| \leq \sup_{\partial \Omega} |Du| + C,
\]
where $C$ is a positive constant depending only on $\lambda$, $\operatorname{diam}(\Omega)$, $|a_{ij}, b_i|_{C^1(\bar{\Omega})}$, $M = |u|_{L^\infty(\Omega)}$ and $|f|_{C^1(\bar{\Omega}\times[-M,M])}$.
\end{theorem}

\begin{proof}
Set $L \equiv a_{ij}D_{ij} + b_iD_i$. We calculate $L(|Du|^2)$ first. Note
\[
D_i(|Du|^2)=2D_kuD_{ki}u,
\]
and
\begin{equation}
D_{ij}(|Du|^2)=2D_{ki}uD_{kj}u+2D_kuD_{kij}u.
\tag{2}
\end{equation}
Differentiating (1) with respect to $x_k$, multiplying by $D_ku$ and summing over $k$, we have by (2)
\[
a_{ij}D_{ij}(|Du|^2)+b_iD_i(|Du|^2)=2a_{ij}D_{ki}uD_{kj}u
\]
\[
- 2D_ka_{ij}D_kuD_{ij}u-2D_kb_iD_kuD_iu+2D_zf|Du|^2+2D_kfD_ku.
\]
The ellipticity assumption implies
\[
\sum_{i,j,k} a_{ij}D_{ki}uD_{kj}u \geq \lambda |D^2u|^2.
\]
By the Cauchy inequality, we have
\[
L(|Du|^2) \geq \lambda |D^2u|^2 - C|Du|^2 - C,
\]
where $C$ is a positive constant depending only on $\lambda$, the $C^1$-norms of $a_{ij}$ and $b_i$ in $\bar{\Omega}$ and the $C^1$-norm of $f$ in $\bar{\Omega}\times[-M,M]$. We need to add another term $u^2$ to control $|Du|^2$ in the right hand side. By the ellipticity assumption again, we have
\[
L(u^2)=2a_{ij}D_iuD_ju+2u\{a_{ij}D_{ij}u+b_iD_i u\}
\]
\[
\geq 2\lambda |Du|^2+2uf.
\]
Therefore we obtain
\[
L(|Du|^2+\alpha u^2)\geq \lambda |D^2u|^2+(2\lambda\alpha-C)|Du|^2-C
\]
\[
\geq \lambda |D^2u|^2+|Du|^2-C,
\]
if we choose $\alpha>0$ large, with $C$ depending in addition on $M$. In order to control the constant term, we consider another function $e^{\beta x_1}$ for $\beta>0$. Hence, we get
\[
L(|Du|^2+\alpha u^2+e^{\beta x_1})\geq \lambda |D^2u|^2+|Du|^2+\{\beta^2a_{11}e^{\beta x_1}+\beta b_1e^{\beta x_1}-C\}.
\]
If we put the domain $\Omega \subset \{x_1>0\}$, then $e^{\beta x_1}\geq 1$ for any $x \in \Omega$. By choosing $\beta$ large, we make the last term positive. Therefore, by taking large $\alpha,\beta$ depending only on $\lambda$, $\operatorname{diam}(\Omega)$, $|a_{ij}|_{C^1(\bar{\Omega})}$, $|b_i|_{C^1(\bar{\Omega})}$, $M = |u|_{L^\infty(\Omega)}$, $|f|_{C^1(\bar{\Omega}\times[-M,M])}$ and setting
\[
w = |Du|^2+\alpha |u|^2+e^{\beta x_1},
\]
we obtain
\[
Lw \geq 0 \quad \text{in } \Omega.
\]
% Source page 0046
By the maximum principle, we have
\[
\sup_{\Omega} w \leq \sup_{\partial\Omega} w.
\]
This finishes the proof.
\end{proof}

Similarly, we discuss the interior gradient estimate. In this case, we just require the bound of \(\sup_{\Omega}|u|\).

\begin{proposition}\label{hanlin:maxprin-prop-2-28}\dcref{Proposition 2.28}\group{ch02-maximum-principles}\source{han-lin-lecture-notes:page-0046}
\textit{Suppose \(u \in C^3(\Omega)\) satisfies}
\[
a_{ij}(x)D_{ij}u + b_i(x)D_i u = f(x,u) \quad \text{in } \Omega,
\]
\textit{for \(a_{ij}, b_i \in C^1(\overline{\Omega})\) and \(f \in C^1(\overline{\Omega}\times \mathbb R)\). Then for any compact subset \(\Omega' \Subset \Omega\)}
\[
\sup_{\Omega'} |Du| \leq C,
\]
\textit{where \(C\) is a positive constant depending on \(\lambda\), \(\diam(\Omega)\), \(\dist(\Omega', \partial\Omega)\), \(|a_{ij}|_{C^1(\overline{\Omega})}\), \(|b_i|_{C^1(\Omega)}\), \(M = |u|_{L^\infty(\Omega)}\) and \(|f|_{C^1(\Omega\times[-M,M])}\).}
\end{proposition}

\begin{proof}
We take a cut-off function \(\gamma \in C_0^\infty(\Omega)\) with \(\gamma \geq 0\) and consider an auxiliary function of the form
\[
w = \gamma |Du|^2 + \alpha |u|^2 + e^{\beta x_1}.
\]
Set \(v = \gamma |Du|^2\). For \(L = a_{ij}D_{ij} + b_iD_i\), we then have
\[
Lv = (L\gamma)|Du|^2 + \gamma L(|Du|^2) + 2a_{ij}D_i\gamma D_j|Du|^2.
\]
Recall in the proof of Theorem 2.27
\[
L(|Du|^2) \geq \lambda |D^2u|^2 - C|Du|^2 - C.
\]
Hence we have
\[
Lv \geq \lambda\gamma|D^2u|^2 + 4a_{ij}D_i\gamma D_k u D_{kj}u - C|Du|^2 + (L\gamma)|Du|^2 - C.
\]
The Cauchy inequality implies for any \(\varepsilon > 0\)
\[
\lvert 4a_{ij}D_k u D_i\gamma D_{kj}u\rvert \leq \varepsilon |D\gamma|^2|D^2u|^2 + c(\varepsilon)|Du|^2.
\]
For the cut-off function \(\gamma\), we require
\[
|D\gamma|^2 \leq C\gamma \quad \text{in } \Omega.
\]
Therefore we have by taking \(\varepsilon > 0\) small
\[
Lv \geq \lambda\gamma|D^2u|^2 \left(1-\varepsilon \frac{|D\gamma|^2}{\gamma}\right) - C|Du|^2 - C
\]
\[
\geq \frac{1}{2}\lambda\gamma|D^2u|^2 - C|Du|^2 - C.
\]
Now we may proceed as before.

In the rest of this section, we use barrier functions to derive boundary gradient estimates. We need to assume that the domain \(\Omega\) satisfies the uniform exterior sphere property.
\end{proof}

% Source page 0047
\section*{2.4. GRADIENT ESTIMATES}

\begin{proposition}\label{hanlin:maxprin-prop-2-29}\dcref{Proposition 2.29}\group{ch02-maximum-principles}\source{han-lin-lecture-notes:page-0047}
\textit{Suppose $u \in C^2(\Omega) \cap C(\overline{\Omega})$ satisfies}
\[
a_{ij}(x)D_{ij}u + b_i(x)D_i u = f(x,u) \quad \text{in } \Omega,
\]
\textit{for $a_{ij}, b_i \in C(\overline{\Omega})$ and $f \in C(\overline{\Omega}\times \mathbb{R})$. Then}
\[
|u(x)-u(x_0)| \leq C|x-x_0| \quad \text{for any $x \in \Omega$ and $x_0 \in \partial\Omega$.}
\]
\textit{where $C$ is a positive constant depending only on $\lambda,\Omega, |a_{ij},b_i|_{L^\infty(\Omega)}, M = |u|_{L^\infty(\Omega)}, |f|_{L^\infty(\Omega\times[-M,M])}$ and $|\varphi|_{C^2(\overline{\Omega})}$ for some $\varphi \in C^2(\overline{\Omega})$ with $\varphi = u$ on $\partial\Omega$.}
\end{proposition}

\begin{proof}
For simplicity, we assume $u = 0$ on $\partial\Omega$. As before, we set $L = a_{ij}D_{ij} + b_iD_i$. Then we have
\[
L(\pm u) = \pm f \geq -F \quad \text{in } \Omega,
\]
where we denote $F = \sup_{\Omega}|f(\cdot,u)|$. Now we fix $x_0 \in \partial\Omega$ and construct a function $w$ such that
\[
Lw \leq -F \text{ in } \Omega,\quad w(x_0)=0,\quad w|_{\partial\Omega}\geq 0.
\]
Then by the maximum principle we have
\[
-w \leq u \leq w \quad \text{in } \Omega.
\]
Taking normal derivatives at $x_0$, we obtain
\[
\left|\frac{\partial u}{\partial n}(x_0)\right| \leq \frac{\partial w}{\partial n}(x_0).
\]
So we need to bound $\dfrac{\partial w}{\partial n}(x_0)$ independently of $x_0$.

Consider an exterior ball $B_R(y)$ with $\overline{B_R(y)} \cap \overline{\Omega} = \{x_0\}$. Define $d(x)$ as the distance from $x$ to $\partial B_R(y)$. Then we have
\[
0 < d(x) < D \equiv \diam(\Omega) \quad \text{for any $x \in \Omega$.}
\]
In fact, $d(x)=|x-y|-R$ for any $x \in \Omega$. Consider $w=\psi(d)$ for some function $\psi$ defined in $[0,\infty)$. Then we need
\[
\psi(0)=0 \qquad (\Rightarrow w(x_0)=0),
\]
\[
\psi(d)>0 \text{ for } d>0 \qquad (\Rightarrow w|_{\partial\Omega}\geq 0),
\]
\[
\psi'(0) \text{ is controlled.}
\]
From the first two inequalities, it is natural to require that $\psi'(d)>0$. Note
\[
Lw = \psi''a_{ij}D_i d D_j d + \psi'a_{ij}D_{ij}d + \psi'b_iD_i d.
\]
A direct calculation yields
\[
D_i d(x)=\frac{x_i-y_i}{|x-y|},
\]
\[
D_{ij}d(x)=\frac{\delta_{ij}}{|x-y|}-\frac{(x_i-y_i)(x_j-y_j)}{|x-y|^3},
\]
which imply $|Dd|=1$ and, with $\Lambda=\sup|a_{ij}|$,
\[
a_{ij}D_{ij}d=\frac{a_{ii}}{|x-y|}-\frac{a_{ij}}{|x-y|}D_i d D_j d
\]
\[
\leq \frac{n\Lambda}{|x-y|}-\frac{\lambda}{|x-y|}=\frac{n\Lambda-\lambda}{|x-y|}\leq \frac{n\Lambda-\lambda}{R}.
\]
\end{proof}

% Source page 0048
Therefore we obtain by the ellipticity
\[
Lw \leq \psi'' a_{ij}D_iD_jd + \psi'\left(\frac{n\Lambda - \lambda}{R} + \lvert b\rvert\right)
\]
\[
\leq \lambda\psi'' + \psi'\left(\frac{n\Lambda - \lambda}{R} + \lvert b\rvert\right),
\]
if we require \(\psi'' < 0\). Hence in order to have \(Lw \leq -F\), we need
\[
\lambda\psi'' + \psi'\left(\frac{n\Lambda - \lambda}{R} + \lvert b\rvert\right) + F \leq 0.
\]

To this end, we study the equation for some positive constants \(a\) and \(b\)
\[
\psi'' + a\psi' + b = 0,
\]
whose solution is given by
\[
\psi(d) = -\frac{b}{a}d + \frac{C_1}{a} - \frac{C_2}{a}e^{-ad},
\]
for some constants \(C_1\) and \(C_2\). For \(\psi(0) = 0\), we need \(C_1 = C_2\). Hence we have for
some constant \(C\)
\[
\psi(d) = -\frac{b}{a}d + \frac{C}{a}(1 - e^{-ad}),
\]
which implies
\[
\psi'(d) = Ce^{-ad} - \frac{b}{a} = e^{-ad}\left(C - \frac{b}{a}e^{ad}\right),
\]
\[
\psi''(d) = -Cae^{-ad}.
\]

In order to have \(\psi'(d) > 0\), we need \(C \geq \frac{b}{a}e^{aD}\). Since \(\psi'(d) > 0\) for \(d > 0\), so
\(\psi(d) > \psi(0) = 0\) for any \(d > 0\). Therefore we take
\[
\psi(d) = -\frac{b}{a}d + \frac{b}{a^2}e^{aD}(1 - e^{-ad})
\]
\[
= \frac{b}{a}\left\{\frac{1}{a}e^{aD}(1 - e^{-ad}) - d\right\}.
\]

Such a \(\psi\) satisfies all the requirements we imposed.

\subsection*{2.5. Alexandroff Maximum Principle}

Suppose \(\Omega\) is a bounded domain in \(\mathbf{R}^n\) and consider a second order elliptic
operator \(L\) in \(\Omega\) of the form
\[
L = a_{ij}(x)D_{ij} + b_i(x)D_i + c(x),
\]
where coefficients \(a_{ij}, b_i, c\) are at least continuous in \(\Omega\). The ellipticity means that
the coefficient matrix \(A = (a_{ij})\) is positive definite everywhere in \(\Omega\). We set \(D =\)
\(\det(A)\) and \(D^* = D^{\frac{1}{n}}\) so that \(D^*\) is the geometric mean of the eigenvalues of \(A\).
Throughout this section, we assume
\[
0 < \lambda \leq D^* \leq \Lambda,
\]
where \(\lambda\) and \(\Lambda\) are two positive constants, which denote the minimal and maximal
eigenvalues of \(A\) respectively.

% Source page 0049
\[
\text{2.5. ALEXANDROFF MAXIMUM PRINCIPLE}
\]

Before stating the main theorem, we first introduce the concept of contact sets.
For $u \in C^2(\Omega)$, we define
\[
\Gamma^+ = \{y \in \Omega;\ u(x) \le u(y) + Du(y)\cdot (x-y)\ \text{for any } x \in \Omega\}.
\]
The set $\Gamma^+$ is called the upper contact set of $u$ and the Hessian matrix
$D^2u = (D_{ij}u)$ is nonpositive on $\Gamma^+$. In fact, the upper contact set
can also be defined for a continuous function $u$ by
\[
\Gamma^+ = \{y \in \Omega;\ u(x) \le u(y) + p\cdot (x-y),
\]
\[
\text{for any } x \in \Omega \text{ and some } p = p(y) \in \mathbf{R}^n\}.
\]
Clearly, $u$ is concave if and only if $\Gamma^+ = \Omega$. If $u \in C^1(\Omega)$,
then $p(y) = Du(y)$ and any support hyperplane must then be a tangent plane to
the graph.

Now we consider an equation of the following form
\[
Lu = f \quad \text{in } \Omega,
\]
for some $f \in C(\Omega)$. The following result is referred to as the Alexandroff
maximum principle.

\begin{theorem}\label{hanlin:maxprin-theorem-2-30}\dcref{Theorem 2.30}\group{ch02-maximum-principles}\source{han-lin-lecture-notes:page-0049,han-lin-lecture-notes:page-0050,han-lin-lecture-notes:page-0051}
\textit{Suppose $u \in C(\overline{\Omega}) \cap C^2(\Omega)$
satisfies $Lu \ge f$ in $\Omega$ with $\frac{|b|}{D^*}, \frac{f}{D^*} \in L^n(\Omega)$
and $c \le 0$ in $\Omega$. Then}
\[
\sup_{\Omega} u \le \sup_{\partial \Omega} u^+ + C\left\|\frac{f^-}{D^*}\right\|_{L^n(\Gamma^+)}.
\]
\textit{where $\Gamma^+$ is the upper contact set of $u$ and $C$ is a constant depending only on $n$,
$\operatorname{diam}(\Omega)$ and $\left\|\frac{b}{D^*}\right\|_{L^n(\Gamma^+)}$. In fact, $C$ can be written as}
\[
d \cdot \left\{\exp\!\left(\frac{2^{\,n-2}}{\omega_n^{\,n}}\left(\left\|\frac{b}{D^*}\right\|_{L^n(\Gamma^+)} + 1\right)\right) - 1\right\},
\]
\textit{with $\omega_n$ as the volume of the unit ball in $\mathbf{R}^n$.}
\end{theorem}

\begin{remark}\label{hanlin:maxprin-remark-2-31}\dcref{Remark 2.31}\group{ch02-maximum-principles}\source{han-lin-lecture-notes:page-0049}
The integral domain $\Gamma^+$ can be replaced by
\[
\Gamma^+ \cap \{x \in \Omega;\ u(x) > \sup_{\partial \Omega} u^+\}.
\]
\end{remark}

\begin{remark}\label{hanlin:maxprin-remark-2-32}\dcref{Remark 2.32}\group{ch02-maximum-principles}\source{han-lin-lecture-notes:page-0049}
There is no assumption on the uniform ellipticity. Compare
with Theorem 2.24.
\end{remark}

To prove Theorem 2.30, we need a lemma first.

\begin{lemma}\label{hanlin:maxprin-lemma-2-33}\dcref{Lemma 2.33}\group{ch02-maximum-principles}\source{han-lin-lecture-notes:page-0049,han-lin-lecture-notes:page-0050}
\textit{Suppose $g \in L^1_{\mathrm{loc}}(\mathbf{R}^n)$ is
nonnegative. Then for any $u \in C(\overline{\Omega}) \cap C^2(\Omega)$}
\[
\int_{B_{\widetilde{M}}(0)} g \le \int_{\Gamma^+} g(Du)\,|\det D^2u|,
\]
\textit{where $\Gamma^+$ is the upper contact set of $u$ and $\widetilde{M} =
(\sup_{\Omega} u - \sup_{\partial \Omega} u^+)/d$ with $d = \operatorname{diam}(\Omega)$.}
\end{lemma}

% Source page 0050
\begin{remark}\label{hanlin:maxprin-remark-2-34}\dcref{Remark 2.34}\group{ch02-maximum-principles}\source{han-lin-lecture-notes:page-0050}\uses{hanlin:maxprin-lemma-2-33}
For any positive definite matrix $A=(a_{ij})$, we have
\[
\det(-D^2u)\le \frac{1}{D}\left(\frac{-a_{ij}D_{ij}u}{n}\right)^n
\quad \text{on } \GammaP.
\]
Hence we have another form for Lemma 2.33
\[
\int_{B_{\Mtilde}(0)} g \le \int_{\GammaP} g(Du)\left(\frac{-a_{ij}D_{ij}u}{nD^*}\right)^n.
\]
\end{remark}

\begin{remark}\label{hanlin:maxprin-remark-2-35}\dcref{Remark 2.35}\group{ch02-maximum-principles}\source{han-lin-lecture-notes:page-0050}\uses{hanlin:maxprin-lemma-2-33}
By taking $g=1$, we have
\[
\sup_{\Omega} u \le \sup_{\partial\Omega} u^+ + \frac{d}{\omega_n^{1/n}}
\left(\int_{\GammaP} \left|\det D^2u\right|\right)^{\!1/n}
\]
\[
\le \sup_{\partial\Omega} u^+ + \frac{d}{\omega_n^{1/n}}
\left(\int_{\GammaP}\left(\frac{-a_{ij}D_{ij}u}{nD^*}\right)^n\right)^{\!1/n}.
\]
This is Theorem 2.30 if $b_i\equiv 0$ and $c=0$.
\end{remark}

\begin{proof}[Proof of Lemma 2.33]
Without loss of generality, we assume $u\le 0$ on $\partial\Omega$.
Set $\OmegaP=\{u>0\}$. By the area-formula for $Du$ in $\GammaP\cap \OmegaP\subset \Omega$, we have
\begin{equation}
\int_{Du(\GammaP\cap\OmegaP)} g \le \int_{\GammaP\cap\OmegaP} g(Du)|\det(D^2u)|
\tag{1}
\end{equation}
where $|\det(D^2u)|$ is the Jacobian of the map $Du:\Omega\to\R^n$. In fact we may consider
$\chi_\el = Du-\el I_d:\Omega\to\R^n$. Then $D\chi_\el = D^2u-\el I$, which is negative definite in $\GammaP$.
Hence by the change of variable formula, we have
\[
\int_{\chi_\el(\GammaP\cap\OmegaP)} g=\int_{\GammaP\cap\OmegaP} g(\chi_\el)|\det(D^2u-\el I)|
\]
which implies (1) if we let $\el\to 0$.
Now we claim $B^M_{\widetilde{M}}(0)\subset Du(\GammaP\cap\OmegaP)$, i.e., for any $a\in\R^n$ with $|a|<\widetilde{M}$ there
exists an $x\in \GammaP\cap\OmegaP$ such that $a=Du(x)$. We may assume $u$ attains its maximum
$m>0$ at $0\in\Omega$, i.e.,
\[
u(0)=m=\sup_{\Omega}u.
\]
Consider an affine function for $|a|<m/d(\equiv \widetilde{M})$
\[
L(x)=m+a\cdot x.
\]
Then $L(x)>0$ for any $x\in\Omega$ and $L(0)=m$. Since u attains its maximum at $0$, then
$Du(0)=0$. Hence there exists an $x_1$ close to $0$ such that $u(x_1)>L(x_1)>0$. Note
that $u\le 0<L$ on $\partial\Omega$. Hence there exists an $\tilde{x}\in\Omega$ such that $Du(\tilde{x})=DL(\tilde{x})=a$.
Now we may translate vertically the plane $y=L(x)$ to the highest such position,
i.e., the whole surface $y=u(x)$ lies below the plane. Clearly at such a point, the
function $u$ is positive.
\hfill$\square$
\end{proof}

\begin{proof}[Proof of Theorem 2.30]\uses{hanlin:maxprin-lemma-2-33}
We should choose $g$ appropriately in order to apply
Lemma 2.33. Note $(-a_{ij}D_{ij}u)^n\le |b|^n|Du|^n$ in $\Omega$ if $f\equiv 0$ and $c\equiv 0$. This
suggests that we should take $g(p)=|p|^{-n}$. However, such a function is not locally

% Source page 0051
integrable (at the origin). Hence, we will choose $g(p) = (|p|^n + \mu^n)^{-1}$ and then let
$\mu \to 0^+$.

First we have by the Cauchy inequality
\[
-a_{ij}D_{ij}u \leq b_iD_i u + cu - f
\]
\[
\leq b_iD_i u - f \qquad \text{in } \Omega^+ = \{x;u(x)>0\}
\]
\[
\leq |b| \cdot |Du| + f^-
\]
\[
\leq \left(|b|^n + \frac{(f^-)^n}{\mu^n}\right)^{\frac{1}{n}} \cdot (|Du|^n + \mu^n)^{\frac{1}{n}} \cdot (1+1)^{\frac{n-2}{n}},
\]
and in particular
\[
(-a_{ij}D_{ij}u)^n \leq \left(|b|^n + \left(\frac{f^-}{\mu}\right)^n\right)(|Du|^n + \mu^n)\cdot 2^{n-2}.
\]

Now we choose
\[
g(p) = \frac{1}{|p|^n + \mu^n}.
\]

By Lemma 2.33, we have
\[
\int_{B_{\tilde M}(0)} g \leq \frac{2^{n-2}}{n^n} \int_{\Gamma^+\cap\Omega^+} \frac{|b|^n + \mu^{-n}(f^-)^n}{D}.
\]

We evaluate the integral in the left hand side in the following way
\[
\int_{B_{\tilde M}(0)} g = \omega_n \int_0^{\tilde M} \frac{r^{n-1}}{r^n + \mu^n}dr = \frac{\omega_n}{n}\log \frac{\tilde M^n + \mu^n}{\mu^n} = \frac{\omega_n}{n}\log\left(\frac{\tilde M^n}{\mu^n} + 1\right).
\]

Therefore we obtain
\[
\tilde M^n \leq \mu^n \left\{\exp\left\{\frac{2^{n-2}}{\omega_n n^n}\left\|\frac{b}{D^*}\right\|_{L^n(\Gamma^+\cap\Omega^+)}^n + \mu^{-n}\left\|\frac{f^-}{D^*}\right\|_{L^n(\Gamma^+\cap\Omega^+)}^n\right\} - 1\right\}.
\]

If $f \neq 0$, we choose $\mu = \left\|\frac{f^-}{D^*}\right\|_{L^n(\Gamma^+\cap\Omega^+)}$. If $f = 0$, we may choose any $\mu > 0$ and
then let $\mu \to 0$.
\hfill$\square$
\end{proof}

In the following, we use Theorem 2.30 and Lemma 2.33 to derive a priori
estimates for solutions of quasilinear equations and fully nonlinear equations. In
the next result, we do not assume the uniform ellipticity.

\begin{theorem}\label{hanlin:maxprin-theorem-2-36}\dcref{Theorem 2.36}\group{ch02-maximum-principles}\source{han-lin-lecture-notes:page-0051,han-lin-lecture-notes:page-0052}
\textit{Suppose $u \in C(\overline{\Omega}) \cap C^2(\Omega)$ satisfies}
\[
Qu = a_{ij}(x,u,Du)D_{ij}u + b(x,u,Du) = 0 \qquad \text{in } \Omega,
\]
\textit{where $a_{ij} \in C(\Omega \times \mathbb{R} \times \mathbb{R}^n)$ satisfies}
\[
a_{ij}(x,z,p)\xi_i\xi_j > 0 \qquad \text{for any $(x,z,p)\in \Omega \times \mathbb{R} \times \mathbb{R}^n$ and $\xi \in \mathbb{R}^n$.}
\]
\textit{Suppose there exist non-negative functions $g \in L^n_{\mathrm{loc}}(\mathbb{R}^n)$ and $h \in L^n(\Omega)$ such that}
\[
\frac{|b(x,z,p)|}{nD^*} \leq \frac{h(x)}{g(p)} \qquad \text{for any $(x,z,p)\in \Omega \times \mathbb{R} \times \mathbb{R}^n$,}
\]
\textit{and}
\[
\int_\Omega h^n(x)\,dx < \int_{\mathbb{R}^n} g^n(p)\,dp \equiv g_\infty.
\]
\textit{Then}
\[
\sup_{\Omega} |u| \le \sup_{\partial \Omega} |u| + C \diam(\Omega),
\]
\textit{where $C$ is a positive constant depending only on $g$ and $h$.}
\end{theorem}

% Source page 0052

\begin{example}\label{hanlin:maxprin-example-2-37}\dcref{Example 2.37}\group{ch02-maximum-principles}\source{han-lin-lecture-notes:page-0052}
The prescribed mean curvature equation is given by
\[
(1 + |Du|^2)\Delta u - D_i u D_j u D_{ij}u = nH(x)(1 + |Du|^2)^{\frac{3}{2}},
\]
for some $H \in C(\Omega)$. We have
\[
a_{ij}(x,z,p) = (1 + |p|^2)\delta_{ij} - p_i p_j,
\]
\[
b(x,z,p) = -nH(x)(1 + |p|^2)^{\frac{3}{2}}.
\]
This implies $D = (1 + |p|^2)^{n-1}$ and
\[
\frac{|b(x,z,p)|}{nD^*} \le \frac{|H(x)|(1 + |p|^2)^{\frac{3}{2}}}{(1 + |p|^2)^{\frac{n-1}{n}}}
= |H(x)|(1 + |p|^2)^{\frac{n+2}{2n}},
\]
and in particular
\[
g_{\infty} = \int_{\mathbb{R}^n} g^n(p)\,dp
= \int_{\mathbb{R}^n} \frac{dp}{(1 + |p|^2)^{\frac{n+2}{2}}}
= \omega_n.
\]
\end{example}

\begin{corollary}\label{hanlin:maxprin-cor-2-38}\dcref{Corollary 2.38}\group{ch02-maximum-principles}\source{han-lin-lecture-notes:page-0052}
Suppose $u \in C(\bar{\Omega}) \cap C^2(\Omega)$ satisfies
\[
(1 + |Du|^2)\Delta u - D_i u D_j u D_{ij}u = nH(x)(1 + |Du|^2)^{\frac{3}{2}}
\quad \text{in } \Omega,
\]
for some $H \in C(\Omega)$ with
\[
H_0 = \int_{\Omega} |H(x)|^n\,dx < \omega_n.
\]
Then
\[
\sup_{\Omega} |u| \le \sup_{\partial \Omega} |u| + C \diam(\Omega),
\]
where $C$ is a positive constant depending only on $n$ and $H_0$.
\end{corollary}

\begin{proof}[Proof of Theorem 2.36]\uses{hanlin:maxprin-lemma-2-33}
We only prove for subsolutions. Assume $Qu \ge 0$
in $\Omega$. Then we have
\[
-a_{ij}D_{ij}u \le b \quad \text{in } \Omega.
\]
Note that $(D_{ij}u)$ is nonpositive in $\Gamma^+$. Hence $-a_{ij}D_{ij}u \ge 0$,
which implies $b(x,u,Du) \ge 0$ in $\Gamma^+$. Then
\[
\frac{b(x,z,Du)}{nD^*} \le \frac{h(x)}{g(Du)}
\quad \text{in } \Gamma^+ \cap \Omega^+.
\]
We apply Lemma 2.33 to $g^n$ and get
\[
\int_{B_{\tilde M}(0)} g^n \le
\int_{\Gamma^+\cap\Omega^+} g^n(Du)\left(\frac{-a_{ij}D_{ij}u}{nD^*}\right)^n
\le \int_{\Gamma^+\cap\Omega^+} g^n(Du)\left(\frac{b}{nD^*}\right)^n
\]
\[
\le \int_{\Gamma^+\cap\Omega^+} h^n \le \int_{\Omega} h^n \left(< \int_{\mathbb{R}^n} g^n\right).
\]
Therefore there exists a positive constant $C$, depending only on $g$ and $h$, such
that $\tilde M \le C$. This implies
\[
\sup_{\Omega} u \le \sup_{\partial \Omega} u^+ + C \diam(\Omega).
\]
This finishes the proof.
\end{proof}

% Source page 0053

Next we discuss Monge-Ampère equations.

\begin{theorem}\label{hanlin:maxprin-theorem-2-39}\dcref{Theorem 2.39}\group{ch02-maximum-principles}\source{han-lin-lecture-notes:page-0053}
Suppose $u \in C(\bar{\Omega}) \cap C^2(\Omega)$ satisfies
\[
\det(D^2u)=f(x,u,Du)\qquad \text{in }\Omega,
\]
for some $f \in C(\Omega\times\mathbb{R}\times\mathbb{R}^n)$. Suppose there exist nonnegative functions $g \in L^1_{\mathrm{loc}}(\mathbb{R}^n)$ and $h \in L^1(\Omega)$ such that
\[
\abs{f(x,z,p)} \le \frac{h(x)}{g(p)} \qquad \text{for any } (x,z,p) \in \Omega \times \mathbb{R} \times \mathbb{R}^n,
\]
and
\[
\int_{\Omega} h(x)\,dx < \int_{\mathbb{R}^n} g(p)\,dp \equiv g_\infty.
\]
Then
\[
\sup_{\Omega}\abs{u} \le \sup_{\partial\Omega}\abs{u} + C\,\diam(\Omega),
\]
where $C$ is a positive constant depending only on $g$ and $h$.
\end{theorem}

The proof is similar to that of Theorem 2.36. There are two special cases. The first case is given by $f=f(x)$. We may take $g=1$ and hence $g_\infty=\infty$.

\begin{corollary}\label{hanlin:maxprin-cor-2-40}\dcref{Corollary 2.40}\group{ch02-maximum-principles}\source{han-lin-lecture-notes:page-0053}
Let $u \in C(\bar{\Omega}) \cap C^2(\Omega)$ satisfy
\[
\det(D^2u)=f(x)\qquad \text{in }\Omega,
\]
for some $f \in C(\bar{\Omega})$. Then
\[
\sup_{\Omega}\abs{u} \le \sup_{\partial\Omega}\abs{u} + \frac{\diam(\Omega)}{\omega_n^{1/n}}\left(\int_{\Omega} \abs{f}^n\right)^{1/n}.
\]
\end{corollary}

The second case is the prescribed Gauss curvature equations.

\begin{corollary}\label{hanlin:maxprin-cor-2-41}\dcref{Corollary 2.41}\group{ch02-maximum-principles}\source{han-lin-lecture-notes:page-0053}
Let $u \in C(\bar{\Omega}) \cap C^2(\Omega)$ satisfy
\[
\det(D^2u)=K(x)(1+\abs{Du}^2)^{\frac{n+2}{2}}\qquad \text{in }\Omega,
\]
for some $K \in C(\bar{\Omega})$ with
\[
K_0 \equiv \int_{\Omega} \abs{K(x)} < \omega_n.
\]
Then
\[
\sup_{\Omega}\abs{u} \le \sup_{\partial\Omega}\abs{u} + C\diam(\Omega),
\]
where $C$ is a positive constant depending only on $n$ and $K_0$.
\end{corollary}

We finish this section by proving a maximum principle in domains with small volumes, which is due to Varadham.

Consider
\[
Lu \equiv a_{ij}D_{ij}u+b_iD_iu+cu \qquad \text{in }\Omega,
\]
where $(a_{ij})$ is positive definite pointwisely in $\Omega$ and
\[
\abs{b_i}+\abs{c} \le \Lambda,\qquad \det(a_{ij})\ge \lambda,
\]
for some positive constants $\lambda$ and $\Lambda$.

% Source page 0054
\begin{theorem}\label{hanlin:maxprin-theorem-2-42}\dcref{Theorem 2.42}\group{ch02-maximum-principles}\source{han-lin-lecture-notes:page-0054}
Suppose $u \in C(\overline{\Omega}) \cap C^2(\Omega)$ satisfies $Lu \ge 0$ in $\Omega$ with $u \le 0$ on $\partial \Omega$. Assume $\operatorname{diam}(\Omega) \le d$. Then there exists a positive constant $\delta$, depending only on $n$, $\lambda$, $\Lambda$ and $d$, such that if $|\Omega| \le \delta$ then $u \le 0$ in $\Omega$.
\end{theorem}

\begin{proof}
If $c < 0$, then $u \le 0$ by Theorem 2.30. In general, we write $c = c^+ - c^-$. Then
\[
a_{ij} D_{ij}u + b_i D_i u - c^-u \ge -c^+u (\equiv f).
\]
By Theorem 2.30, we have
\[
\sup_{\Omega} u \le C \|c^+u\|_{L^n(\Omega)} \le C \|c^+\|_{L^\infty(\Omega)} |\Omega|^{1/n} \cdot \sup_{\Omega} u \le \frac{1}{2}\sup_{\Omega} u,
\]
if $|\Omega|$ is small, where $C$ is a positive constant, depending only on $n$, $\lambda$, $\Lambda$ and $d$.
Hence we get $u \le 0$ in $\Omega$.
\end{proof}

Compare this with Proposition 2.15, the maximum principle for narrow domains.

\section*{2.6. The Moving Plane Method}

In this section, we revisit the moving plane method to discuss the symmetry of solutions.

\begin{theorem}\label{hanlin:maxprin-theorem-2-43}\dcref{Theorem 2.43}\group{ch02-maximum-principles}\source{han-lin-lecture-notes:page-0054,han-lin-lecture-notes:page-0055}
Suppose $u \in C(\overline{B}_1) \cap C^2(B_1)$ is a positive solution of
\[
\Delta u + f(u) = 0 \qquad \text{in } B_1
\]
\[
u = 0 \qquad \text{on } \partial B_1,
\]
where $f$ is locally Lipschitz in $\mathbb{R}$. Then $u$ is radially symmetric in $B_1$ and $\frac{\partial u}{\partial r}(x) < 0$ for $x \ne 0$.
\end{theorem}

We already proved Theorem 2.43 if $u$ is $C^2$ up to the boundary. Here we give a method which does not depend on the smoothness of domains nor the smoothness of solutions up to the boundary.

\begin{lemma}\label{hanlin:maxprin-lemma-2-44}\dcref{Lemma 2.44}\group{ch02-maximum-principles}\source{han-lin-lecture-notes:page-0055}
Suppose that the bounded domain $\Omega$ is convex in $x_1$ direction and symmetric with respect to the hyperplane $\{x_1 = 0\}$ and that $u \in C(\overline{\Omega}) \cap C^2(\Omega)$ is a positive solution of
\[
\Delta u + f(u) = 0 \qquad \text{in } \Omega
\]
\[
u = 0 \qquad \text{on } \partial \Omega,
\]
where $f$ is locally Lipschitz in $\mathbb{R}$. Then $u$ is symmetric with respect to $x_1$ and $D_{x_1}u(x) < 0$ for any $x \in \Omega$ with $x_1 > 0$.
\end{lemma}

\begin{proof}
Write $x = (x_1,y) \in \Omega$ for $y \in \mathbb{R}^{n-1}$. We will prove
\[
(1)\qquad u(x_1,y) < u(x_1^*,y) \text{ for any } x_1 > 0 \text{ and } x_1^* < x_1 \text{ with } x_1^* + x_1 > 0.
\]
Then by letting $x_1^* \to -x_1$, we have $u(x_1,y) \leq u(-x_1,y)$ for any $x_1$. By changing the direction $x_1 \to -x_1$, we get the symmetry.

% Source page 0055
\[
\text{2.6. THE MOVING PLANE METHOD}\qquad 49
\]

Let \(a=\sup x_1\) for \((x_1,y)\in\Omega\). For \(0<\lambda<a\), define
\[
\Sigma_\lambda=\{x\in\Omega;x_1>\lambda\},
\]
\[
T_\lambda=\{x_1=\lambda\},
\]
\[
\Sigma'_\lambda=\text{the reflection of }\Sigma_\lambda\text{ with respect to }T_\lambda,
\]
\[
x_\lambda=(2\lambda-x_1,x_2,\cdots x_n)\text{ for }x=(x_1,x_2,\cdots x_n).
\]

In \(\Sigma_\lambda\), we define
\[
w_\lambda(x)=u(x)-u(x_\lambda)\qquad \text{for }x\in\Sigma_\lambda.
\]

Then we have by the mean value theorem
\[
\Delta w_\lambda+c(x,\lambda)w_\lambda=0\qquad \text{in }\Sigma_\lambda,
\]
\[
w_\lambda\leq 0\text{ and }w_\lambda\neq 0\qquad \text{on }\partial\Sigma_\lambda,
\]
where \(c(x,\lambda)\) is a bounded function in \(\Sigma_\lambda\).

We need to show \(w_\lambda<0\) in \(\Sigma_\lambda\) for any \(\lambda\in(0,a)\). This implies in particular
that \(w_\lambda\) assumes along \(\partial\Sigma_\lambda\cap\Omega\) its maximum in \(\Sigma_\lambda\). By Theorem 2.6, Hopf Lemma,
we have for any such \(\lambda\in(0,a)\)
\[
D_{x_1}w_\lambda\bigg|_{x_1=\lambda}=2D_{x_1}u\bigg|_{x_1=\lambda}<0.
\]

For any \(\lambda\) close to \(a\), we have \(w_\lambda<0\) by Proposition 2.15, the maximum
principle for narrow domains, or Theorem 2.42, the maximum principle for domains
with small volumes. Let \((\lambda_0,a)\) be the largest interval of values of \(\lambda\) such that
\(w_\lambda<0\) in \(\Sigma_\lambda\). We want to show \(\lambda_0=0\). If \(\lambda_0>0\), by the continuity, we have \(w_{\lambda_0}\leq 0\)
in \(\Sigma_{\lambda_0}\) and \(w_{\lambda_0}\neq 0\) on \(\partial\Sigma_{\lambda_0}\). Then Theorem 2.8, the strong maximum principle,
implies \(w_{\lambda_0}<0\) in \(\Sigma_{\lambda_0}\). We will show that for any small \(\varepsilon>0\)
\[
w_{\lambda_0-\varepsilon}<0\qquad \text{in }\Sigma_{\lambda_0-\varepsilon}.
\]

Fix a \(\delta>0\) to be determined. Let \(K\) be a closed subset in \(\Sigma_{\lambda_0}\) such that \(|\Sigma_{\lambda_0}\setminus K|<\)
\(\delta/2\). The fact \(w_{\lambda_0}<0\) in \(\Sigma_{\lambda_0}\) implies
\[
w_{\lambda_0}(x)\leq -\eta<0\qquad \text{for any }x\in K.
\]

By the continuity, we have
\[
w_{\lambda_0-\varepsilon}<0\qquad \text{in }K.
\]

For any \(\varepsilon>0\) small, \(|\Sigma_{\lambda_0-\varepsilon}\setminus K|<\delta\). We choose \(\delta\) in such a way that we apply
Theorem 2.42 to \(w_{\lambda_0-\varepsilon}\) in \(\Sigma_{\lambda_0-\varepsilon}\setminus K\). Hence we get
\[
w_{\lambda_0-\varepsilon}(x)\leq 0\qquad \text{in }\Sigma_{\lambda_0-\varepsilon}\setminus K,
\]
and then by Theorem 2.12
\[
w_{\lambda_0-\varepsilon}(x)<0\qquad \text{in }\Sigma_{\lambda_0-\varepsilon}\setminus K.
\]

Therefore we obtain for any small \(\varepsilon>0\)
\[
w_{\lambda_0-\varepsilon}(x)<0\qquad \text{in }\Sigma_{\lambda_0-\varepsilon}.
\]

This contradicts the choice of \(\lambda_0\) and hence finishes the proof.
\end{proof}
